\chapter{Einschränkungen, Komposition und Kürzen}
\label{b06guide:ch:composition}

\section{Bereiche anpassen}

\begin{definition}[Einschränkung, Komposition und Identität]
Für $f\colon A\to B$ und $U\subseteq A$ ist
$f|_U=\{(x,f(x))\mid x\in U\}$ die Einschränkung auf $U$.
Für $g\colon B\to C$ wird die Komposition durch
$(g\circ f)(x)=g(f(x))$ für $x\in A$ bezeichnet.
Ferner sei $\operatorname{id}_A\colon A\to A$, $x\mapsto x$,
die Identität auf $A$.
\end{definition}

\begin{theorem}[Injektivität von Einschränkungen]
\label{b06guide:thm:restriction}
Seien $f\colon A\rightarrowtail B$, $U\subseteq A$ und
$V\subseteq B$. Dann sind
\[
 f|_U\colon U\rightarrowtail B,
 \qquad
 f|_{f^{-1}[V]}\colon f^{-1}[V]\rightarrowtail V
\]
injektiv. Für jede Menge $D$ mit $f[U]\subseteq D$ ist
auch $f|_U\colon U\rightarrowtail D$ injektiv.
\end{theorem}

\begin{theorem}[Unabhängigkeit von der zulässigen Zielmenge]
\label{b06guide:thm:codomain}
Seien $f\colon A\to B$ und $D$ eine Menge mit $f[A]\subseteq D$.
Dann gilt für denselben Graphen
\[
 f\colon A\rightarrowtail B
 \quad\Longleftrightarrow\quad
 f\colon A\rightarrowtail D.
\]
Insbesondere erhalten sowohl das Vergrößern der Zielmenge als
auch ihre Beschränkung auf eine Menge, die alle Werte enthält,
die Injektivität.
\end{theorem}

\section{Kompositionen erkennen und kürzen}

\begin{theorem}[Genaues Injektivitätskriterium einer Komposition]
\label{b06guide:thm:composition}
Für $f\colon A\to B$ und $g\colon B\to C$ gilt
\[
 g\circ f\text{ injektiv}
 \quad\Longleftrightarrow\quad
 \left\{
 \begin{array}{l}
 f\text{ ist injektiv},\\
 g|_{f[A]}\text{ ist injektiv}.
 \end{array}\right.
\]
Daraus ergeben sich insbesondere:
\begin{enumerate}
\item Die Komposition zweier injektiver Funktionen ist injektiv.
\item Ist $g\circ f$ injektiv, so ist $f$ injektiv.
\item Ist $g$ injektiv, so ist $g\circ f$ genau dann injektiv,
      wenn $f$ injektiv ist.
\end{enumerate}
Die Injektivität von $g$ auf ganz $B$ folgt im allgemeinen
nicht aus der Injektivität der Komposition.
\end{theorem}

\begin{theorem}[Linkes Kürzen als Charakterisierung]
\label{b06guide:thm:cancellation}
Für $f\colon A\to B$ sind äquivalent:
\begin{enumerate}
\item $f$ ist injektiv.
\item Für jede Menge $T$ und alle Funktionen $u,v\colon T\to A$ gilt
\[
 f\circ u=f\circ v\quad\Longrightarrow\quad u=v.
\]
\item Für eine fest gewählte Einermenge $T$ und alle
      $u,v\colon T\to A$ gilt dieselbe Kürzungsregel.
\end{enumerate}
Anstelle der Einermenge darf jede fest gewählte nichtleere Menge
$T$ verwendet werden. Die leere Testmenge genügt nicht.
\end{theorem}

% Die zusammengehörige Definition und ihre beiden Sätze gemeinsam beginnen.
\Needspace{28\baselineskip}
\section{Linksinversen und eindeutige Faktorisierung}

\begin{definition}[Linksinverse]
Für $f\colon A\to B$ heißt eine Funktion $r\colon B\to A$
eine \emph{Linksinverse} von $f$, wenn
\[
 r\circ f=\operatorname{id}_A.
\]
Es wird allein diese Gleichheit verlangt.
\end{definition}

\begin{theorem}[Charakterisierung durch eine Linksinverse]
\label{b06guide:thm:leftinverse}
Für jede Funktion $f\colon A\to B$ gilt
\[
 \begin{gathered}
 \exists r\colon B\to A\quad
          r\circ f=\operatorname{id}_A\\
 \Longleftrightarrow\\
 f\text{ ist injektiv und }(A\ne\varnothing\text{ oder }B=\varnothing).
 \end{gathered}
\]
Für nichtleeres $A$ ist die Existenz einer Linksinversen also
gleichwertig mit der Injektivität. Die leere Funktion
$\varnothing\to B$ ist auch für $B\ne\varnothing$ injektiv,
hat dann aber keine Linksinverse.
\end{theorem}

\begin{theorem}[Eindeutige Faktorisierung durch eine Injektion]
\label{b06guide:thm:factorization}
Seien $f\colon A\rightarrowtail B$ und $g\colon C\to B$.
Es existiert genau dann eine Funktion $h\colon C\to A$ mit
\[
 f\circ h=g,
\]
wenn $g[C]\subseteq f[A]$. In diesem Fall ist $h$ eindeutig
bestimmt. Außerdem gilt
\[
 h\text{ injektiv}\quad\Longleftrightarrow\quad g\text{ injektiv}.
\]
Die Aussage schließt leere Mengen ein.
\end{theorem}